%! AP Statistics — Unit 3, Lesson 3.3 — Guided Notes KEY
\documentclass[10pt]{article}
\usepackage{apstats-article}
\usepackage{apstats-key}

\begin{document}

\pageheader{Unit 3, Lesson 3.3}{Guided Notes — KEY}
\namedateperiod

\vspace{0.12in}

\begin{objectivebox}
\small By the end of this lesson, I will be able to\dots\\[4pt]
\ans{Describe and distinguish SRS, stratified, cluster, systematic random sampling, and census.}\\[1pt]
\ans{Identify the sampling method used in a given scenario.}\\[1pt]
\ans{Explain when each method is most appropriate.}
\end{objectivebox}

\vspace{0.1in}

\begin{vocabbox}
\termblanklong{Simple random sample (SRS)} \ans{Every group of size $n$ has an equal chance of being selected; individuals chosen using a random device.}
\termblanklong{Stratified random sample} \ans{Divide population into homogeneous strata; take an SRS from each stratum.}
\termblanklong{Cluster sample} \ans{Divide population into heterogeneous clusters; randomly select clusters and survey all individuals in them.}
\termblanklong{Systematic random sample} \ans{Choose a random start; select every $k$th individual from a list.}
\termblanklong{Census} \ans{Measure every individual in the population.}
\end{vocabbox}

\vspace{0.1in}

\begin{notesbox}{Sampling Methods Reference}
\small
\renewcommand{\arraystretch}{2.0}
\begin{tabularx}{\linewidth}{@{} >{\bfseries\color{navy}}p{0.24\linewidth}
    >{\raggedright\arraybackslash}p{0.38\linewidth}
    >{\raggedright\arraybackslash}X @{}}
\textbf{Method} & \textbf{How to do it} & \textbf{Best when \ldots} \\
\hline
Simple Random Sample &
Number all; randomly select $n$. Every group of size $n$ has \ans{equal} chance. &
The population is \ans{relatively homogeneous} \\
Stratified Random Sample &
Divide into \ans{homogeneous} groups called \ans{strata}. SRS within each stratum. &
You want to ensure representation of \ans{subgroups} \\
Cluster Sample &
Divide into \ans{clusters} (heterogeneous). Randomly select clusters; survey
\ans{everyone} in selected clusters. &
Individuals are \ans{spread out} geographically \\
Systematic Random Sample &
Choose a \ans{random} start; select every \ans{$k$}th individual. &
You have a complete \ans{ordered} list \\
Census &
Measure \ans{every individual} in the population. &
Population is \ans{small} and manageable \\
\end{tabularx}
\end{notesbox}

\vspace{0.1in}

\begin{notesbox}{Stratified vs. Cluster --- KEY Distinction}
\small
\begin{tabularx}{\linewidth}{@{} >{\bfseries\color{navy}}p{0.3\linewidth} X @{}}
Strata are \ldots & \ans{homogeneous} (same type within each group) \\
Clusters are \ldots & \ans{heterogeneous} (mixed types within each group) \\
In stratified: survey & \ans{some} from \ans{every} stratum \\
In cluster: survey & \ans{everyone} from \ans{randomly selected} clusters \\
\end{tabularx}
\end{notesbox}

\vspace{0.1in}

\begin{practicebox}
\small
\begin{enumerate}[leftmargin=1.5em, itemsep=10pt]
  \item A researcher numbers all 2{,}000 employees and uses a random number table to select 100. \ans{Simple random sample (SRS)}

  \item A school divides students into 4 grade levels, then randomly selects 25 students from each grade. \ans{Stratified random sample}

  \item A city randomly selects 6 of its 60 neighborhoods; surveys every household in the selected neighborhoods. \ans{Cluster sample}

  \item A manager randomly picks a start between 1--10, then surveys every 10th customer. \ans{Systematic random sample}
\end{enumerate}
\end{practicebox}

\end{document}
